% Options for packages loaded elsewhere
\PassOptionsToPackage{unicode}{hyperref}
\PassOptionsToPackage{hyphens}{url}
\documentclass[
]{article}
\usepackage{xcolor}
\usepackage[margin=1in]{geometry}
\usepackage{amsmath,amssymb}
\setcounter{secnumdepth}{-\maxdimen} % remove section numbering
\usepackage{iftex}
\ifPDFTeX
  \usepackage[T1]{fontenc}
  \usepackage[utf8]{inputenc}
  \usepackage{textcomp} % provide euro and other symbols
\else % if luatex or xetex
  \usepackage{unicode-math} % this also loads fontspec
  \defaultfontfeatures{Scale=MatchLowercase}
  \defaultfontfeatures[\rmfamily]{Ligatures=TeX,Scale=1}
\fi
\usepackage{lmodern}
\ifPDFTeX\else
  % xetex/luatex font selection
\fi
% Use upquote if available, for straight quotes in verbatim environments
\IfFileExists{upquote.sty}{\usepackage{upquote}}{}
\IfFileExists{microtype.sty}{% use microtype if available
  \usepackage[]{microtype}
  \UseMicrotypeSet[protrusion]{basicmath} % disable protrusion for tt fonts
}{}
\makeatletter
\@ifundefined{KOMAClassName}{% if non-KOMA class
  \IfFileExists{parskip.sty}{%
    \usepackage{parskip}
  }{% else
    \setlength{\parindent}{0pt}
    \setlength{\parskip}{6pt plus 2pt minus 1pt}}
}{% if KOMA class
  \KOMAoptions{parskip=half}}
\makeatother
\usepackage{color}
\usepackage{fancyvrb}
\newcommand{\VerbBar}{|}
\newcommand{\VERB}{\Verb[commandchars=\\\{\}]}
\DefineVerbatimEnvironment{Highlighting}{Verbatim}{commandchars=\\\{\}}
% Add ',fontsize=\small' for more characters per line
\usepackage{framed}
\definecolor{shadecolor}{RGB}{248,248,248}
\newenvironment{Shaded}{\begin{snugshade}}{\end{snugshade}}
\newcommand{\AlertTok}[1]{\textcolor[rgb]{0.94,0.16,0.16}{#1}}
\newcommand{\AnnotationTok}[1]{\textcolor[rgb]{0.56,0.35,0.01}{\textbf{\textit{#1}}}}
\newcommand{\AttributeTok}[1]{\textcolor[rgb]{0.13,0.29,0.53}{#1}}
\newcommand{\BaseNTok}[1]{\textcolor[rgb]{0.00,0.00,0.81}{#1}}
\newcommand{\BuiltInTok}[1]{#1}
\newcommand{\CharTok}[1]{\textcolor[rgb]{0.31,0.60,0.02}{#1}}
\newcommand{\CommentTok}[1]{\textcolor[rgb]{0.56,0.35,0.01}{\textit{#1}}}
\newcommand{\CommentVarTok}[1]{\textcolor[rgb]{0.56,0.35,0.01}{\textbf{\textit{#1}}}}
\newcommand{\ConstantTok}[1]{\textcolor[rgb]{0.56,0.35,0.01}{#1}}
\newcommand{\ControlFlowTok}[1]{\textcolor[rgb]{0.13,0.29,0.53}{\textbf{#1}}}
\newcommand{\DataTypeTok}[1]{\textcolor[rgb]{0.13,0.29,0.53}{#1}}
\newcommand{\DecValTok}[1]{\textcolor[rgb]{0.00,0.00,0.81}{#1}}
\newcommand{\DocumentationTok}[1]{\textcolor[rgb]{0.56,0.35,0.01}{\textbf{\textit{#1}}}}
\newcommand{\ErrorTok}[1]{\textcolor[rgb]{0.64,0.00,0.00}{\textbf{#1}}}
\newcommand{\ExtensionTok}[1]{#1}
\newcommand{\FloatTok}[1]{\textcolor[rgb]{0.00,0.00,0.81}{#1}}
\newcommand{\FunctionTok}[1]{\textcolor[rgb]{0.13,0.29,0.53}{\textbf{#1}}}
\newcommand{\ImportTok}[1]{#1}
\newcommand{\InformationTok}[1]{\textcolor[rgb]{0.56,0.35,0.01}{\textbf{\textit{#1}}}}
\newcommand{\KeywordTok}[1]{\textcolor[rgb]{0.13,0.29,0.53}{\textbf{#1}}}
\newcommand{\NormalTok}[1]{#1}
\newcommand{\OperatorTok}[1]{\textcolor[rgb]{0.81,0.36,0.00}{\textbf{#1}}}
\newcommand{\OtherTok}[1]{\textcolor[rgb]{0.56,0.35,0.01}{#1}}
\newcommand{\PreprocessorTok}[1]{\textcolor[rgb]{0.56,0.35,0.01}{\textit{#1}}}
\newcommand{\RegionMarkerTok}[1]{#1}
\newcommand{\SpecialCharTok}[1]{\textcolor[rgb]{0.81,0.36,0.00}{\textbf{#1}}}
\newcommand{\SpecialStringTok}[1]{\textcolor[rgb]{0.31,0.60,0.02}{#1}}
\newcommand{\StringTok}[1]{\textcolor[rgb]{0.31,0.60,0.02}{#1}}
\newcommand{\VariableTok}[1]{\textcolor[rgb]{0.00,0.00,0.00}{#1}}
\newcommand{\VerbatimStringTok}[1]{\textcolor[rgb]{0.31,0.60,0.02}{#1}}
\newcommand{\WarningTok}[1]{\textcolor[rgb]{0.56,0.35,0.01}{\textbf{\textit{#1}}}}
\usepackage{graphicx}
\makeatletter
\newsavebox\pandoc@box
\newcommand*\pandocbounded[1]{% scales image to fit in text height/width
  \sbox\pandoc@box{#1}%
  \Gscale@div\@tempa{\textheight}{\dimexpr\ht\pandoc@box+\dp\pandoc@box\relax}%
  \Gscale@div\@tempb{\linewidth}{\wd\pandoc@box}%
  \ifdim\@tempb\p@<\@tempa\p@\let\@tempa\@tempb\fi% select the smaller of both
  \ifdim\@tempa\p@<\p@\scalebox{\@tempa}{\usebox\pandoc@box}%
  \else\usebox{\pandoc@box}%
  \fi%
}
% Set default figure placement to htbp
\def\fps@figure{htbp}
\makeatother
\setlength{\emergencystretch}{3em} % prevent overfull lines
\providecommand{\tightlist}{%
  \setlength{\itemsep}{0pt}\setlength{\parskip}{0pt}}
\usepackage{bookmark}
\IfFileExists{xurl.sty}{\usepackage{xurl}}{} % add URL line breaks if available
\urlstyle{same}
\hypersetup{
  pdftitle={brms Basics},
  hidelinks,
  pdfcreator={LaTeX via pandoc}}

\title{brms Basics}
\author{}
\date{\vspace{-2.5em}2026-04-17}

\begin{document}
\maketitle

\subsection{Introduction}\label{introduction}

\texttt{brms} is the most widely used high-level interface between R and
Stan. It accepts the familiar \texttt{lme4} formula syntax for fixed and
random effects, automatically translates the model to Stan, compiles it,
runs the No-U-Turn Sampler (NUTS), and returns a posterior together with
diagnostic and predictive helpers. This means you can fit Bayesian
generalised linear, mixed, multinomial, ordinal, censored,
distributional, or non-linear models without writing Stan by hand, while
still being able to inspect or override the generated code when the
situation demands it.

\subsection{Prerequisites}\label{prerequisites}

Comfort with classical regression, an understanding of priors and
posteriors, and a working \texttt{rstan} or \texttt{cmdstanr} toolchain
(a C++ compiler is required on first use). Familiarity with
\texttt{lme4} formulas accelerates the learning curve dramatically
because almost every random-effects expression carries over unchanged.

\subsection{Theory}\label{theory}

A \texttt{brms} model has three pieces: a \textbf{formula} (mean
structure plus optional distributional terms), a \textbf{family}
(likelihood and link), and \textbf{priors} (per parameter class). Random
effects use \texttt{(slope\ \textbar{}\ group)} syntax, with
\texttt{\textbar{}\textbar{}} for uncorrelated terms. Distributional
models, set up through
\texttt{bf(y\ \textasciitilde{}\ ...,\ sigma\ \textasciitilde{}\ ...)},
let any parameter of the family depend on predictors. Priors are
attached by class: \texttt{b} for population-level slopes,
\texttt{Intercept} for the centred intercept, \texttt{sd} for
random-effect standard deviations, \texttt{sigma} for residual scale,
and \texttt{cor} for correlation matrices via the LKJ distribution.
Defaults are weakly informative but should be reviewed for any analysis
intended for publication.

\subsection{Assumptions}\label{assumptions}

Inherits the assumptions of the chosen family --- for \texttt{gaussian}
that means linearity in the linear predictor, additivity, and Normal
residuals after partial pooling. Convergence diagnostics
(\(\hat R < 1.01\), no divergent transitions, sensible effective sample
size) are also a precondition for trusting any posterior summary.

\subsection{R Implementation}\label{r-implementation}

\begin{Shaded}
\begin{Highlighting}[]
\FunctionTok{library}\NormalTok{(brms)}

\FunctionTok{data}\NormalTok{(sleepstudy, }\AttributeTok{package =} \StringTok{"lme4"}\NormalTok{)}

\NormalTok{fit }\OtherTok{\textless{}{-}} \FunctionTok{brm}\NormalTok{(Reaction }\SpecialCharTok{\textasciitilde{}}\NormalTok{ Days }\SpecialCharTok{+}\NormalTok{ (Days }\SpecialCharTok{|}\NormalTok{ Subject),}
           \AttributeTok{data =}\NormalTok{ sleepstudy,}
           \AttributeTok{family =}\NormalTok{ gaussian,}
           \AttributeTok{prior =} \FunctionTok{prior}\NormalTok{(}\FunctionTok{normal}\NormalTok{(}\DecValTok{0}\NormalTok{, }\DecValTok{30}\NormalTok{), }\AttributeTok{class =} \StringTok{"b"}\NormalTok{) }\SpecialCharTok{+}
                   \FunctionTok{prior}\NormalTok{(}\FunctionTok{exponential}\NormalTok{(}\FloatTok{0.1}\NormalTok{), }\AttributeTok{class =} \StringTok{"sigma"}\NormalTok{) }\SpecialCharTok{+}
                   \FunctionTok{prior}\NormalTok{(}\FunctionTok{exponential}\NormalTok{(}\FloatTok{0.1}\NormalTok{), }\AttributeTok{class =} \StringTok{"sd"}\NormalTok{) }\SpecialCharTok{+}
                   \FunctionTok{prior}\NormalTok{(}\FunctionTok{lkj}\NormalTok{(}\DecValTok{2}\NormalTok{), }\AttributeTok{class =} \StringTok{"cor"}\NormalTok{),}
           \AttributeTok{chains =} \DecValTok{4}\NormalTok{, }\AttributeTok{iter =} \DecValTok{2000}\NormalTok{, }\AttributeTok{refresh =} \DecValTok{0}\NormalTok{)}

\FunctionTok{summary}\NormalTok{(fit)}
\FunctionTok{pp\_check}\NormalTok{(fit)}
\end{Highlighting}
\end{Shaded}

\subsection{Output \& Results}\label{output-results}

\texttt{summary()} reports posterior means, standard errors, 95 \%
credible intervals, \(\hat R\), bulk effective sample size, and tail
effective sample size for every population-level coefficient,
group-level standard deviation, and the residual scale.
\texttt{pp\_check()} overlays simulated datasets on the observed
outcome; if the empirical density falls inside the cloud of replicated
densities, the model captures the marginal distribution adequately.

\subsection{Interpretation}\label{interpretation}

A reporting sentence: ``A Bayesian linear mixed model (\texttt{brms}
2.21, four chains, 2,000 iterations) estimated reaction time to increase
by 10.5 ms per day of sleep deprivation (95 \% CrI 7.4 to 13.6), with
substantial between-subject heterogeneity in both intercept
(\(SD = 26\)) and slope (\(SD = 6.5\)); \(\hat R = 1.00\) for all
parameters.'' Always pair the point estimate with the credible interval
and a convergence statistic.

\subsection{Practical Tips}\label{practical-tips}

\begin{itemize}
\tightlist
\item
  Run \texttt{get\_prior()} on your formula before fitting; it lists
  every prior class so you can override defaults instead of inheriting
  them silently.
\item
  Inspect \texttt{stancode(fit)} and \texttt{standata(fit)} when
  something behaves oddly; the generated code is readable and reveals
  exactly how \texttt{brms} parameterised the model.
\item
  \texttt{conditional\_effects()} produces marginal effect plots with
  credible bands and is the fastest way to communicate non-linear or
  interaction effects to non-Bayesian audiences.
\item
  Use \texttt{add\_criterion(fit,\ "loo")} to attach LOO-CV for later
  model comparison; storing it on the object avoids re-fitting.
\item
  For very large datasets, switch to \texttt{backend\ =\ "cmdstanr"} and
  \texttt{threads\_per\_chain\ =\ 2} to exploit within-chain
  parallelism.
\item
  When defaults yield divergent transitions, increase
  \texttt{adapt\_delta} to 0.95 or 0.99 before reaching for stronger
  priors.
\end{itemize}

\subsection{R Packages Used}\label{r-packages-used}

\texttt{brms} for model specification and compilation, \texttt{rstan} or
\texttt{cmdstanr} as the sampling backend, \texttt{bayesplot} for
diagnostic plots used internally by \texttt{pp\_check()}, and
\texttt{posterior} for handling draw arrays.

\subsection{For Reviewers}\label{for-reviewers}

\textbf{What to look for in a paper using this method.}

\begin{itemize}
\tightlist
\item
  \textbf{Common misapplications.}
\item
  \textbf{Diagnostics that should be reported but often aren't.}
\item
  \textbf{Red flags in tables and figures.}
\item
  \textbf{What to verify.}
\item
  \textbf{What an adequate Methods paragraph must contain.}
\end{itemize}

\subsection{See also --- labs in this
chapter}\label{see-also-labs-in-this-chapter}

\begin{itemize}
\tightlist
\item
  \href{labs/lab-bayesian-thinking-control-vs-decision-error.Rmd}{Bayesian
  thinking; α control vs decision error}
\item
  \href{labs/lab-brms-stan-loo-hierarchical-models.Rmd}{brms / Stan,
  LOO, hierarchical models}
\end{itemize}

\end{document}
